\chapter{Introduction générale}

\section*{Contexte général}
\addcontentsline{toc}{section}{Contexte général}

Les réseaux cellulaires 4G/5G supportent aujourd'hui un volume de
trafic considérable et en croissance constante, porté par la
multiplication des usages mobiles (données Internet, messagerie,
voix). Cette densification s'accompagne d'une complexité croissante
dans la supervision du réseau~: les opérateurs doivent être en mesure
de détecter rapidement tout comportement de trafic sortant de la
norme, qu'il s'agisse d'une panne, d'une saturation localisée ou d'un
événement exceptionnel affectant une zone géographique donnée. La
détection d'anomalies dans les séries temporelles de trafic mobile
constitue ainsi un enjeu opérationnel important pour la qualité de
service.

\section*{Problématique}
\addcontentsline{toc}{section}{Problématique}

La difficulté principale réside dans la distinction entre une
variation normale du trafic (heures de pointe, week-ends, saisonnalité)
et un comportement réellement atypique. À l'échelle d'une métropole,
cette analyse doit en outre être menée sur un grand nombre de zones
géographiques et sur plusieurs canaux de communication (Internet, SMS,
appels), dont les dynamiques diffèrent sensiblement. Une méthode de
détection doit donc être suffisamment sensible pour repérer les
journées atypiques, tout en restant robuste au bruit statistique
inhérent à des données réelles à large échelle.

\section*{Objectifs}
\addcontentsline{toc}{section}{Objectifs}

Ce travail vise à concevoir et évaluer une méthode de détection
d'anomalies dans le trafic de réseaux cellulaires 4G/5G, appliquée à
des données réelles de trafic mobile de la ville de Milan. Les
objectifs poursuivis sont les suivants~:
\begin{itemize}
  \item encoder le profil horaire de trafic de chaque zone sous une
        forme compacte et comparable d'un jour à l'autre~;
  \item détecter, pour chaque zone, les journées présentant un
        comportement statistiquement atypique~;
  \item appliquer cette détection indépendamment sur les trois canaux
        disponibles (Internet, SMS, appels) et évaluer l'apport d'une
        confirmation multi-canal~;
  \item valider les anomalies détectées par recoupement avec un
        calendrier d'événements connus (jours fériés, matchs à
        domicile à San Siro).
\end{itemize}

\section*{Démarche adoptée}
\addcontentsline{toc}{section}{Démarche adoptée}

L'approche retenue s'inspire du framework académique DiNATrAX
(Maudoux \& Boumerdassi, IEEE ICC 2024) : le trafic horaire de chaque
zone est encodé en une signature journalière de 24 lettres, chaque
lettre représentant un niveau de trafic exprimé en multiple de la
médiane horaire de la zone. Cette signature est comparée à un profil
de référence par la distance d'édition de Damerau-Levenshtein, puis
convertie en score z permettant d'isoler les jours anormaux. Le détail
de cette méthode est présenté au chapitre~\ref{chap:methodologie}.

Le jeu de données utilisé est résumé dans le tableau~\ref{tab:dataset}.

\begin{table}[htbp]
  \centering
  \caption{Jeu de données utilisé}
  \label{tab:dataset}
  \begin{tabular}{ll}
    \toprule
    \textbf{Caractéristique} & \textbf{Valeur} \\
    \midrule
    Source & Telecom Italia Big Data Challenge \\
    Dépôt & Harvard Dataverse \\
    Licence & ODbL \\
    Zone couverte & Ville de Milan (grille géographique) \\
    Période couverte & 1\textsuperscript{er} novembre 2013 -- 1\textsuperscript{er} janvier 2014 \\
    Nombre de fichiers & 62 (un par jour) \\
    Volume par fichier & \textasciitilde{}4 millions de lignes (\textasciitilde{}250 Mo) \\
    Canaux disponibles & Internet, SMS, appels \\
    \bottomrule
  \end{tabular}
\end{table}

\section*{Organisation du rapport}
\addcontentsline{toc}{section}{Organisation du rapport}

La suite de ce rapport est organisée comme indiqué dans le
tableau~\ref{tab:plan}.

\begin{table}[htbp]
  \centering
  \caption{Plan du rapport}
  \label{tab:plan}
  \begin{tabular}{ll}
    \toprule
    \textbf{Chapitre} & \textbf{Contenu} \\
    \midrule
    Contexte et état de l'art & Réseaux 4G/5G, détection d'anomalies, framework DiNATrAX \\
    Méthodologie & Encodage en signatures, distance d'édition, score z, pipeline multi-canal \\
    Résultats & Anomalies détectées, validation calendaire, tableau de bord 3D \\
    Conclusion & Bilan, limites et perspectives \\
    \bottomrule
  \end{tabular}
\end{table}
